\section{Intervals and trivial cofibrations}\label{sec:B}

We keep the distinction between weak saturation (pushouts, transfinite
composites and retracts) and saturation by two out of three.  This distinction
is essential in cote 105/007 and 105/012--020.  Write
$\mathscr E=\widehat A=\operatorname{Fun}(A^{\mathrm{op}},\mathrm{Set})$
for a presheaf category on a small category, $1$ for its terminal object and
$C$ for its monomorphisms.  A sectioned interval here means exactly a pair
$(I,e)$ with $e:1\to I$ in $\mathscr E$; it need not have two chosen endpoints.
For $i:K\hookrightarrow L$ in $\mathscr E$ put
\[
 h(i,I,e): (L\times\{e\})\cup_{K\times\{e\}}(K\times I)
              \longrightarrow L\times I .
\]
The braces denote the image of the section, not a choice of a vertex in each
fibre.  All these maps are monomorphisms.  For a specified bunch $\mathcal I$
of sectioned intervals, let $J_{\mathcal I}$ be this class of maps,
$F=J_{\mathcal I}^{\perp}$, $TF=C^{\perp}$, and
$TC=\operatorname{wsat}(J_{\mathcal I})$.  Whenever the small object argument
applies, $TC={}^{\perp}F$.  The negative results below hold with either
interpretation of $TC$.

\begin{claim}{B.1}{Proved}{What hypotheses (1)--(5) actually prove}\label{B1}
Assume $W$ contains identities, satisfies two out of three and
Loc~(1c), that $L_A\times X\to X$ belongs to $W$ for
every $X$, that $I\times X\to X$ belongs to $W$ for every $(I,e)\in\mathcal I$
and every $X$, and that $C\cap W$ is closed under pushouts, filtered colimits
and retracts.  Then $TC\subseteq C\cap W$ and $TF\subseteq W$.  If
$\mathcal I$ is a set, both $(C,TF)$ and $(TC,F)$ are cofibrantly generated
weak factorization systems.  No accessibility hypothesis on $W$ is needed
for these factorizations.
\end{claim}
\begin{proof}
The section $X\to X\times I$ belongs to $W$ by two out of three.  Forming its
pushout along $K\hookrightarrow L$ gives a $W$-map
\[
 L\longrightarrow (L\times\{e\})\cup(K\times I).
\]
Its composite with $h(i,I,e)$ is the section $L\to L\times I$ and belongs
to $W$.  Thus $h(i,I,e)\in C\cap W$.  Successor compositions and limit stages
of a transfinite composite remain in $W$, the latter by filtered-colimit
closure; retracts remain in $W$ by hypothesis.  This proves the first
inclusion.

For the second inclusion let $p:E\to B$ belong to $TF$.  Lifting against
$\varnothing\hookrightarrow B$ gives $s:B\to E$ with $ps=1_B$.
The two disjoint sections of the Lawvere interval give a monomorphism
$E\amalg E\to E\times L_A$.  The lifting property for $p$ gives a map
$H:E\times L_A\to E$ over $B$ whose endpoint restrictions are $sp$ and
$1_E$.  The projection and its endpoints belong to $W$.  Two out of three
then gives $H\in W$ and $sp\in W$.  Loc~(1c) applied to $p,s$, since
$ps=1_B$, gives $p\in W$.

For the size assertion choose an infinite regular cardinal $\kappa$ larger
than the cardinality of the arrow set of $A$.  Up to isomorphism the inclusions
$K\hookrightarrow L$ with fewer than $\kappa$ elements form a set $M$.
Every monomorphism is a transfinite composite of pushouts of elements of
$M$: well-order the elements of its target outside the source, adjoining
successively the subpresheaf generated by the next element.  This
subpresheaf and its intersection with the preceding stage have fewer than
$\kappa$ elements.  This also proves $\operatorname{wsat}(M)=C$.

For fixed $(I,e)$, the pushout-product operation $i\mapsto h(i,I,e)$
preserves pushouts, transfinite composites and retracts.  This can be checked
objectwise in sets, where products preserve colimits; equivalently it follows
from the pushout-product/exponential adjunction.  Consequently
\[
 J=\{h(m,I,e):m\in M,\ (I,e)\in\mathcal I\}
\]
has $J^{\perp}=J_{\mathcal I}^{\perp}$ and the same weak saturation.
It is a set.  Choosing a regular cardinal larger than the sizes of the domains
and codomains of its members makes these objects small for transfinite
composites.  The usual small object construction therefore gives the asserted
weak factorization systems.  If $\mathcal I$ is a proper class, this argument
requires a set of maps with the same right orthogonal; its existence must be
supplied separately.
\end{proof}

\begin{claim}{B.2}{Proved}{The exact additional recognition hypothesis}\label{B2}
Under the hypotheses of B.1 and the existence of both factorizations, the
following are equivalent:
\begin{enumerate}
\item $TC=C\cap W$;
\item every map in $F\cap W$ has the right lifting property against every
monomorphism, that is, $F\cap W\subseteq TF$;
\item $W=TF\circ TC$.
\end{enumerate}
Thus smallness supplies factorizations, whereas condition (2) is the exact
remaining recognition hypothesis.  It cannot be replaced by accessibility
of $W$, by using all universally $W$-aspherical sectioned intervals, or by
requiring a separating interval; B.4 refutes all three substitutions at once.
\end{claim}
\begin{proof}
Suppose (1) and let $p:E\to B$ be in $F\cap W$.  Factor it as
$E\xrightarrow{i}Z\xrightarrow{q}B$ with $i\in C$ and $q\in TF$.
Since $q,p\in W$, we have $i\in C\cap W=TC$.
A lift in the square with top $1_E$, left $i$, right $p$ and bottom $q$
exhibits $p$ as a retract of $q$ in the arrow category.  Hence $p\in TF$.
Conversely, assuming (2), factor $i:K\to L$ in $C\cap W$ as
$K\xrightarrow{j}Z\xrightarrow{p}L$ with $j\in TC$, $p\in F$.
B.1 and two out of three give $p\in F\cap W$, hence $p\in TF$.
A lift against $i$ exhibits $i$ as a retract of $j$, so $i\in TC$.
The same factorization applied to any $w\in W$ proves (3).
Finally (3) applied to $i\in C\cap W$, followed by the same retract square,
proves (1).  Every map and every retract used here is in $\mathscr E$ or its
arrow category, before localization.
\end{proof}

\begin{claim}{B.3}{Refuted}{The total localizer and the empty object}\label{B3}
Even for the accessible total localizer on simplicial sets, all sectioned
intervals together do not generate all monomorphisms.  The map
$\varnothing\hookrightarrow\Delta^0$ is a counterexample.
\end{claim}
\begin{proof}
The class of maps preserving emptiness, meaning that the source is empty
if and only if the target is empty, is closed under pushouts along arbitrary
maps, transfinite composites and retracts.  A generator $h(i,I,e)$ preserves
emptiness: its source contains $L\times\{e\}\cong L$, and if $L$ is empty
so is its target.  Hence every element of $TC$ preserves emptiness, whereas
$\varnothing\hookrightarrow\Delta^0$ does not.  More directly,
$\varnothing\to\Delta^0$ is itself in $F$: every lifting square with a
nonempty source on the left has no top map to $\varnothing$, and a generator
with empty source also has empty target.  It is not in $TF$.
All of (1)--(5) hold for the total class.  The argument uses no size bound and
also separates ${}^{\perp}F$ from $C$.
\end{proof}

\begin{claim}{B.4}{Refuted}{All acyclic intervals still miss homology trivial cofibrations}\label{B4}
Let $A=\Delta$, $\mathscr E=\mathrm{sSet}$, and let $W$ consist of integral
homology equivalences.  This is an accessible localizer, induced by the
accessible fundamental localizer $W_{\mathbb Z}$, and $\Delta$ is a strict
$W_{\mathbb Z}$-test category.  Take \emph{all} pairs $(I,e)$ for which
$I\to\Delta^0$ is an integral homology equivalence.  Then
\[
 TC\subsetneq C\cap W.
\]
The same failure holds for every set or subbunch of these intervals, whether
or not it includes a separating interval.  In particular, the proposed
universal assertion fails in both cases (a) and (b) of cote 105/020; in (b)
one may take $A_0=e$ and $T=\Delta$.
\end{claim}
\begin{proof}
Let $G=S_3$ and let $Z=NG$ be the nerve of its one-object groupoid.  It is a
Kan complex and models $K(S_3,1)$.  We first prove
\[
                 Z\longrightarrow\Delta^0\quad\text{belongs to }F.
\]
Every allowed $I$ is connected and has $H_1(I;\mathbb Z)=0$, so
$\pi_1(I,e)$ is perfect.  The derived series of $S_3$ has orders $6,3,1$.
A homomorphism from a perfect group to a solvable group is trivial: its image
is perfect and is contained successively in every derived subgroup.
It follows that the based mapping space $\operatorname{Map}_*(I,Z)$ has
one component.  Its higher homotopy groups are also zero, since for $n\geq1$
they are represented by maps from the simply connected space
$S^n\wedge I$ to the 1-type $Z$.  These assertions also follow from the
fundamental-groupoid classification of maps to a groupoid nerve, with no
finiteness assumption on $I$.

Evaluation at the chosen section,
\[
            e^*:Z^I\longrightarrow Z,
\]
is a Kan fibration because $e:\Delta^0\hookrightarrow I$ is a cofibration.
Its fibres are the based mapping spaces just computed, so it is a trivial
Kan fibration and has the right lifting property against every simplicial
monomorphism.  By the exponential adjunction, a lift against
$h(i,I,e)$ for $Z\to\Delta^0$ is exactly a lift against $i$ for $e^*$.
This proves the assertion about $F$ for the \emph{entire} class of intervals.
Consequently every element of $TC$, or of ${}^{\perp}F$, must have the
extension property for maps to $Z$.

We now give a finite counterexample to that extension property.  Let $R$
be a bouquet of two circles, based at their common vertex, with generators
$a,b$ of its fundamental group.  Define a graph map $f:R\to R$ on loops by
\[
             f_*(a)=a,\qquad f_*(b)=b[a,b]=bab a^{-1}b^{-1}.
\]
Subdivide the source loops so this is a simplicial map between finite
one-dimensional simplicial sets $X,Y$.  One fully explicit choice is to use
three edges for each target circle, three edges for the source $a$-circle,
and fifteen edges for its $b$-circle, following successively the five letters
$b,a,b,a^{-1},b^{-1}$.  The orientation of each source edge is chosen to map
to the orientation of its target edge.  The common basepoint is the only
identification between the two source circles.  Thus $X$ has 17 vertices
and 18 edges; $Y$ has 5 vertices and 6 edges.

On $H_1(-;\mathbb Z)$ this map is the identity matrix, since the exponent
vectors of the displayed words are $(1,0)$ and $(0,1)$.  It also induces the
identity on $H_0$, and higher homology of both graphs vanishes.
Form the simplicial mapping cylinder
\[
       M_f=Y\amalg_{X\times\{1\}}(X\times\Delta^1),
       \qquad j:X=X\times\{0\}\hookrightarrow M_f.
\]
This pushout is a finite simplicial set and $j$ is a monomorphism.  The
retraction $r:M_f\to Y$, given by $f\operatorname{pr}_X$ on the cylinder
and the identity on $Y$, is a simplicial homotopy equivalence.  Since
$rj=f$, $j$ is an integral homology equivalence.

Choose $g:X\to Z$ whose fundamental-group homomorphism sends
\[
       a\longmapsto(23),\qquad b\longmapsto(12).
\]
Such a map is explicit: label one positively traversed edge in each source
circle with the indicated group element and all its other edges with the
identity.  An extension $M_f\to Z$ would restrict to a map $v:Y\to Z$
and a homotopy between $g$ and $vf$.  Thus the pair $((23),(12))$ would be
simultaneously conjugate to
\[
             (x,\,yxyx^{-1}y^{-1})
                  \quad\text{for some }x,y\in S_3.
\]
It is not.  The attached exhaustive program
\texttt{code/B\_interval\_counterexample.py} enumerates all 36 pairs;
the image has 30 elements, and none of the six simultaneous conjugates of
$((23),(12))$ lies in it.  This can also be checked without computation.
When $x$ is a transposition, $yxyx^{-1}y^{-1}$ is a conjugate of $y$.
If $y$ is the identity or a 3-cycle it cannot be the distinct transposition
$(12)$.  If $y=x$, its value is $x$.  If $y$ is either other transposition,
then $(yx)y(yx)^{-1}=x$ as well.  Hence a transposition $x$ only yields the
transposition $x$ as its second coordinate.  The desired pair is impossible,
and simultaneous conjugation preserves its impossibility.

Therefore $j\notin TC$, although $j\in C\cap W$.  All hypotheses (1)--(5)
hold here: products with an integral-acyclic $I$ preserve integral homology,
pushouts along monomorphisms are controlled by the relative normalized
chain complex, filtered colimits commute with normalized chains and homology,
and homology isomorphisms are retract-closed.  The Lawvere interval is
injective and weakly contractible.  Accessibility and the fundamental-localizer
statement are the integral case of the thesis, Theorem 7.3.2 and
Corollary 7.3.3; the published versions are Theorem 9.4.2 and
Corollary 9.4.3.  The proof of the counterexample itself uses neither an
accessibility assumption nor existence of a factorization for the proper
class of all intervals.
\end{proof}

The conceptual distinction in B.4 has a literature precedent.  Nullification
at all integral-acyclic spaces yields a plus-construction, whereas integral
homology localization inverts more maps.  We opened
\citet[Lemma 3.1, Theorems 3.3 and 4.3]{BerrickCasacuberta1999} and used
Lemma 3.1 only as a comparison; the direct $S_3$ argument above supplies the
needed mapping-space calculation independently.  We claim no priority for
this distinction or for the resulting counterexample.

\begin{claim}{B.5}{Proved}{A single interval can suffice for a specified localizer}\label{B5}
For simplicial sets and ordinary weak homotopy equivalences, the bunch
consisting of $(\Delta^1,0)$ and $(\Delta^1,1)$ gives
$TC=C\cap W_\infty$.  The same single geometric interval with both endpoints
does not work for integral homology equivalences.  Thus cote 105/020 correctly
rejects a fixed separating interval as a solution for arbitrary prescribed
$W$, but it does not rule out a complete interval presentation of a
particular $W$.  The choice of sections matters: the bunch containing only
$(\Delta^1,0)$ fails even for $W_\infty$.
\end{claim}
\begin{proof}
For the two endpoint maps of $\Delta^1$, the weakly saturated class of the
maps $h(i,\Delta^1,\varepsilon)$ is the class generated by the horn
inclusions.  One first reduces $i$ to the boundary inclusions by B.1;
triangulating the prisms $\Delta^n\times\Delta^1$ gives the usual horn
presentation.  This equality is stated in the attached thesis, 3.1.3 and
Proposition 3.1.4 (published 2.1.3--2.1.4).  Its right orthogonal is therefore
the Kan fibrations.  A Kan fibration which is a weak homotopy equivalence
is a trivial fibration, so B.2 applies.  This last recognition is the thesis,
Theorem 3.1.29, or the published Theorem 2.1.42.  B.4 proves the negative
assertion even after adding every other integral-acyclic interval.

For the final assertion, $p:\{1\}\hookrightarrow\Delta^1$ has the right
lifting property against every $h(i,\Delta^1,0)$.  In a lifting problem,
the map $L\times\Delta^1\to\Delta^1=N[1]$ is 1 on
$L\times\{0\}$.  Every vertex $(x,1)$ is the target of the edge from
$(x,0)$, and an edge of $[1]$ with source 1 has target 1.  Thus this map
is 1 on every vertex and hence is constant, since a simplex of $N[1]$
is determined by its ordered vertices.  The lift exists.  But $p$ is a
weak homotopy equivalence and is not right orthogonal to all
monomorphisms.  B.2 therefore disproves equality for this one-section
bunch.  This also explains why the positive assertion explicitly
includes both endpoints.
\end{proof}

\begin{claim}{B.6}{Known}{Accessible localizers admit complete generators, not necessarily interval generators}\label{B6}
For an accessible $A$-localizer $W$, there exists a set $S\subset C\cap W$
with $\operatorname{wsat}(S)=C\cap W$.  Hence adjoining $S$ to the interval
generators repairs the equality.  Accessibility by itself does not assert
that $S$ can consist of the maps $h(i,I,e)$ of cote 105.
\end{claim}
\begin{proof}
The precise source is the attached thesis, Theorem 2.3.4, and the published
Theorem 1.4.3 with Remark 1.4.4.  They construct a cofibrantly generated model
structure with monomorphisms as cofibrations and $W$ as weak equivalences.
By the definition of cofibrant generation its trivial cofibrations have a
set $S$ of generators.  Adding interval generators does not change their
weak saturation, because B.1 places those generators in $C\cap W$.
Conversely, the existence of such a set $S$ makes $W$ accessible as an
$A$-localizer by the same theorem.  The statement concerns a set of arbitrary
monomorphisms, not only a set of objects mapping to $1$.

In the two test-category cases in the prompt, an accessible fundamental
localizer induces this model structure by the thesis, Corollary 5.1.11
(published Corollary 4.2.18).  This supplies complete generators in both
cases.  It does not repair the purely interval-generated proposal, already
refuted in their common special case $\mathscr E=\mathrm{sSet}$.
\end{proof}

\begin{claim}{B.7}{Proved under stated hypotheses}{Two separate and exact tests for a bunch}\label{B7}
Suppose the bunch is a set, contains both endpoints of a separating interval,
and satisfies B.1.  Let $\mathcal V$ be the $A$-localizer associated with its
anodyne class by Cisinski's construction.  Then $TC=C\cap W$ holds if and
only if both of the following hold:
\begin{enumerate}
\item $\mathcal V=W$;
\item for every $Y\in\mathscr E$, the localizer built from the same anodyne
maps over $Y$ is the class of maps over $Y$ whose underlying maps belong to
$\mathcal V$.
\end{enumerate}
Equivalently, the first condition identifies the intended weak equivalences,
and the second asserts completeness of the homotopical structure.  These
conditions are the exact recognition hypothesis of B.2 resolved into its
two independent questions.  They are not consequences of a size bound.
\end{claim}
\begin{proof}
Pushout-product associativity gives
$h(i,I,e)\mathbin{\square}m=h(i\mathbin{\square}m,I,e)$ for every
monomorphism $m$, with the canonical isomorphism of sources and targets
coming from associativity of products and unions.  Since $i\mathbin{\square}m$
is a monomorphism, $TC$ is stable under these pushout-products.  The
separating interval in the bunch therefore makes $TC$ an anodyne class for
an elementary homotopical datum.  Set generation is B.1.

Cisinski's completeness criterion identifies condition (2) with
$TC=C\cap\mathcal V$: this is thesis Proposition 2.2.44, supplemented by
Proposition 2.2.38, and published Propositions 1.3.61 and 1.3.47.
Together with (1) it implies the asserted equality.  Conversely, if
$TC=C\cap W$, then B.2 gives $W=TF\circ TC$.  The resulting model structure
has weak equivalences $W$, and the anodyne construction gives that same
structure; hence $\mathcal V=W$.  The same completeness criterion gives
(2).  In particular, we have not substituted the unproved equality
$\widetilde W=W$ in the proof of itself.
\end{proof}

\subsection{Numbering crosswalk and the historical repairs}

\begin{claim}{B.8}{Known}{The five requested Cisinski comparisons}\label{B8}
The requested numbering does not transfer verbatim between the thesis and
the expanded book.  The precise comparisons are as follows.
\begin{center}\small
\begin{tabular}{p{.16\linewidth}p{.20\linewidth}p{.55\linewidth}}
\hline
Ast\'erisque 308 & Thesis (2002) & Content and qualification\\
\hline
Proposition 3.4.25 & Proposition 4.4.22 & Every $\Delta$-localizer is regular.\\
Theorem 4.2.15 & Theorem 5.1.10 & Test-localizer/fundamental-localizer correspondence, preserving accessibility.\\
Theorem 5.2.15 & Theorem B.2.10 & Accessible fundamental localizers have the left proper Thomason model structure; right properness is characterized.\\
Theorem 6.1.11 & Theorem 5.2.21 and Corollary 5.2.20 & Properness versus generation by aspherical objects; the thesis assumes accessibility at the outset.\\
Theorem 6.1.18 & No exact counterpart located; compare Corollary 5.1.12 & The book proves minimality among \emph{weak} fundamental localizers. The thesis comparison proves minimality among fundamental localizers.\\
\hline
\end{tabular}
\end{center}
The equality problem itself is treated most directly by thesis
Propositions 2.2.38 and 2.2.44 and Theorem 2.3.4, corresponding to book
Propositions 1.3.47 and 1.3.61 and Theorem 1.4.3.  They supply recognition
and existence of complete generators.  B.4 adds the explicit check that
the \emph{maximal} interval bunch in the manuscript does not supply those
generators for an ordinary, accessible homology localizer.

For precision, the correspondence in the second row is the
inclusion-preserving bijection
\[
 \{\text{test }A\text{-localizers }\mathcal V\}
  \ \underset{W\mapsto i_A^{-1}W}{\overset{\mathcal V\mapsto
  \mathcal V_{\mathrm{Cat}}}{\rightleftarrows}}\
 \{\text{fundamental localizers modelable by }A\}.
\]
The third row gives weak equivalences exactly $W$ and cofibrations exactly
the Thomason cofibrations.  In the fourth row the book's five equivalent
conditions are: (i) accessibility and aspherical generation;
(ii) generation by a set of functors $D\to e$;
(iii) being the least fundamental localizer making a set of functors
aspherical; (iv) being the least one modelable by a set of small
categories; and (v) properness.  The thesis lists the counterparts of
(ii)--(v) after assuming accessibility, and its Corollary 5.2.20 separately
states properness for a small family of generators $D\to e$.
The last row says precisely $W_\infty\subseteq W$ for every weak
fundamental localizer $W$, with $W_\infty$ itself a weak fundamental
localizer.  These are statements of the named results, not identifications
of interval-generated anodyne maps with all trivial cofibrations.
\end{claim}
\begin{proof}[Verification of the crosswalk]
We opened both texts and compared the hypotheses and conclusions, not
only the numbers.  The regularity statements identify the homotopy colimit
of representables over the category of elements with the original
simplicial set.  The correspondence statements identify $W$ with the
pullback by the category-of-elements functor of its corresponding
fundamental localizer, and explicitly preserve accessibility.  The Thomason
statements have the same cofibrations, weak equivalences and properness
qualification.  In the aspherical-generation theorem the book packages
accessibility and generation by a set among its equivalent conditions,
whereas thesis Theorem 5.2.21 fixes accessibility first.  Finally, searching
the entire extracted thesis for the weak-localizer terminology and checking
its minimality result found only Corollary 5.1.12; that result has the stronger
relative Theorem A premise through the definition of a fundamental
localizer.  It cannot be cited as the book's weak-localizer theorem.
\end{proof}

\begin{claim}{B.9}{Proved}{What the earlier returns do and do not supply}\label{B9}
BISITE A.27 and A.28 and PATHS D.18 do not prove the missing interval
recognition.  Their conclusions remain valid under the hypotheses stated
there, but those hypotheses already require the necessary exactness or
model-structure input.  The small-set proviso in cote 105/040 and XIX,
Theorem 1.1, is an existence proviso for a factorization argument, not a
proof of the identification $TC=C\cap W$.
\end{claim}
\begin{proof}
In BISITE A.27 the prescribed equality $W=TF\circ TC$ is used to factor a
map $i\in C\cap W$ as $pj$ with $j\in TC$ and $p\in TF$.  The lift
against $i,p$ exhibits $i$ as a retract of $j$; this is exactly the
$(3)\Rightarrow(1)$ proof of B.2.  We rechecked that square and its two
retraction identities.  BISITE A.28 assumes those closed exact structures
on every diagram category before deriving the Kan adjunctions, so it
cannot establish the missing hypothesis for a new interval bunch.
PATHS D.18 invokes the Thomason existence theorem and therefore gives a
model for accessible fundamental localizers; it says nothing identifying
the manuscript's interval-generated class with its trivial cofibrations.
PATHS D.22 explicitly leaves the nonaccessible recognition question open.

The small-set hypothesis of cote 105/040 is precisely what lets the
transfinite construction use a fixed collection of lifting problems with
controlled domains.  B.1 makes that requirement explicit for interval
bunches.  B.4 retains accessibility and refutes the equality, so deleting or
strengthening the small-set hypothesis cannot solve the different
recognition problem.  If $W$ is nonaccessible, B.1 still applies to a set of
intervals, but the existence of a complete set $S$ as in B.6 would imply
accessibility and cannot be assumed.  This is the exact scope in which the
large-cardinal discussion of REVLOC H.1--H.4 is relevant; no large-cardinal
principle is used in any proof in this section.
\end{proof}

\begin{claim}{B.10}{Proved}{On simplicial sets, all intervals generate the aspherical core}\label{B10}
Let $\mathcal W$ be an accessible fundamental localizer, let
$W=\mathcal W_\Delta$ be its localizer on simplicial sets, and use all
sectioned intervals $(I,e)$ with $I\to\Delta^0$ in $W$.
Suppose $\mathcal W$ is not the total localizer.  Write
$\mathcal W_{\mathrm{asph}}$ for the smallest fundamental localizer
containing the collapses of all $\mathcal W$-aspherical categories, and put
$V=(\mathcal W_{\mathrm{asph}})_\Delta$.  Then
\[
 TC_{\mathrm{all}}=C\cap V,
 \qquad F_{\mathrm{all}}=\operatorname{Fib}_V.
\]
A set of the allowed sectioned intervals already gives these equalities.
Consequently, for the bunch of all intervals,
\[
 TC_{\mathrm{all}}=C\cap W
 \quad\Longleftrightarrow\quad
 \mathcal W\text{ is proper}.
\]
Here properness has Cisinski's precise meaning: equivalently,
$\mathcal W$ is accessible and generated by aspherical collapses, or its
associated simplicial model structure is right proper.

For the total localizer the exact answer instead is
\[
 TC_{\mathrm{all}}
 =\{K\hookrightarrow L:K=\varnothing\Longleftrightarrow
                                  L=\varnothing\}.
\]
Thus accessibility alone does not give interval recognition; for a
non-total accessible fundamental localizer on simplicial sets, properness
is the necessary and sufficient additional condition.
\end{claim}
\begin{proof}
The aspherical core and its properties are
\citet[6.1.7, Proposition 6.1.8 and Theorem 6.1.11]{Cisinski2006B}.
It is proper and accessible, is contained in $\mathcal W$, and has exactly
the same aspherical categories.  The simplex-category comparison transfers
this equality to aspherical simplicial sets.  It therefore suffices first
to prove the result for a non-total proper localizer $V$.

\emph{A set of actual sectioned intervals.}
Choose the set $\mathcal A_0$ of aspherical simplicial sets furnished by
\citet[Lemma 6.1.6]{Cisinski2006B}, enlarging it to contain the standard
simplices.  Its defining property is that, for every simplicial set $Y$,
the regular localizer on $\mathrm{sSet}/Y$ generated by maps between
members of $\mathcal A_0$ over $Y$ contains every map between
$V$-aspherical simplicial sets over $Y$.  Such a set is supplied explicitly
there by representatives of the bounded aspherical objects of Lemma 6.1.4.
Since the localizer is non-total, these objects are nonempty.  Each has
vertices, so
\[
 \mathcal I_0=\{(I,e):I\in\mathcal A_0,\ e\in I_0\}
\]
is a set of sectioned intervals; here $I_0$ on the right denotes the set of
vertices of $I$.  In particular both sectioned copies of $\Delta^1$ occur.
If $S^0$ is aspherical, include it in $\mathcal A_0$ as well.  Write
$TC_0,F_0$ for the resulting classes.  B.1 gives both weak factorization
systems and $TC_0\subseteq C\cap V$.  Since both endpoints of $\Delta^1$
are present, B.5 shows that every member of $F_0$ is a Kan fibration.

\emph{The disconnected case.}
If an aspherical simplicial set is disconnected, $S^0$ is a retract of it:
choose vertices in two components and use the component map to send these
components to different vertices of $S^0$.  Hence $S^0$ is aspherical.
The generator $h(i,S^0,e)$ is $L\amalg K\to L\amalg L$ for
$i:K\hookrightarrow L$.  If $p:E\to Y$ is in $F_0$ and $E$ has a vertex,
any lifting square against $i$ extends to one against this generator by
using the chosen vertex constantly on its first copy of $L$.  Thus $p$
has the right lifting property against every monomorphism.
If also $p\in V$ and $Y$ is nonempty, then $E$ cannot be empty: otherwise
$\varnothing\to\Delta^0$ is a retract of $\varnothing\to Y$ and belongs
to $V$, which makes the fundamental localizer total.  When $Y$ is empty,
$p$ is the identity of the empty object.  We obtain
$F_0\cap V\subseteq TF$, and B.2 proves the desired equality in this case.

\emph{Local triviality over the connected aspherical generators.}
We may now assume every $V$-aspherical simplicial set is connected.  Let
$p:E\to Y$ lie in $F_0$, let $v:I\to Y$ be any map with
$I\in\mathcal A_0$, and choose its designated vertex $e$.  Put
$E_I=I\times_YE$ and $Z_e=\{e\}\times_IE_I$.  These are strict pullbacks;
$E_I\to I$ is a Kan fibration and $Z_e$ is Kan.  The lift against
$h(\varnothing\hookrightarrow Z_e,I,e)$ gives a map over $I$
\[
 H:Z_e\times I\longrightarrow E_I,
 \qquad H|_{Z_e\times\{e\}}=1_{Z_e}.
\]
This is a map of Kan fibrations over $I$.  Transport along an edge of $I$
identifies their homotopy fibres, so its equivalence on the fibre at $e$
implies equivalence on every fibre, since $I$ is connected.  Thus $H$ is
an ordinary weak equivalence.  If $Z_e$ is empty, all fibres and $E_I$ are
empty, and the same conclusion holds.  Pulling $H$ back along any
$u:J\to I$ gives the corresponding ordinary weak equivalence over $J$.
For $J\in\mathcal A_0$ the induced map $E_J\to E_I$ is therefore a
$V$-equivalence: under these comparisons it is $Z_e\times u$, and products
preserve $V$-equivalences.

\emph{From the bounded tests to a weak $V$-fibration.}
In the presheaf category $\mathrm{sSet}/Y$, define
\[
 V_p=\{u:U\to U'\text{ over }Y:
          E\times_YU\to E\times_YU'\text{ lies in }V\}.
\]
This is a localizer: pullback preserves colimits and monomorphisms, and
preserves maps having the right lifting property against all monomorphisms.
It is regular.  Indeed $p$ is a Kan fibration, so its pullback functor
preserves ordinary weak equivalences; it also preserves homotopy colimits.
For the latter assertion use the simplicial replacement formula: base
change commutes with its coproducts, simplex factors and realization, and
Kan base change preserves the weak equivalences used in that formula.
The canonical homotopy-colimit comparisons from representables, which
express regularity, consequently remain $V$-equivalences after this base
change.

The preceding paragraph about $H$ puts every map between members of
$\mathcal A_0$ over $Y$ in $V_p$.  Lemma 6.1.6 puts every map between
aspherical simplicial sets over $Y$ in $V_p$.  The criterion of
\citet[Proposition 6.1.1]{Cisinski2006B}, applied to $A=\Delta$, now says
that $p$ is a weak $V$-fibration: pulling back any $V$-equivalence along
$p$ preserves $V$-equivalence.

\emph{A weak equivalence in $F_0$ is a trivial Kan fibration.}
Suppose also $p\in V$, and fix a vertex $y$ of $Y$.  Factor its map in the
$V$-model structure as
\[
 \Delta^0\xrightarrow{a}D\xrightarrow{q}Y,
 \qquad a\in C\cap V,\quad q\in\operatorname{Fib}_V.
\]
Right properness gives $E_D=E\times_YD\to D$ in $V$.  The weak
$V$-fibration property just proved gives
$Z_y=E\times_Y\Delta^0\to E_D$ in $V$.  Since $a\in V$, two out of three
shows $Z_y\to\Delta^0$ is in $V$.  In particular $Z_y$ is nonempty.

The fibre $Z_y\to\Delta^0$ belongs to $F_0$.  For every
$(I,e)\in\mathcal I_0$, its lifting property and the exponential adjunction
make evaluation $Z_y^I\to Z_y$ a trivial Kan fibration.  Hence $Z_y$ is
local with respect to $I\to\Delta^0$ for every $I\in\mathcal A_0$, and
therefore to every map between members of $\mathcal A_0$.
More explicitly, the class of maps $u$ for which
$\operatorname{Map}(u,Z_y)$ is a weak equivalence is a simplicial localizer;
all simplicial localizers are regular by Proposition 3.4.25.  Proposition
6.1.9 therefore makes this class contain $V$.  Applying it to the
aspherical collapse $Z_y\to\Delta^0$ shows that the constant-map inclusion
\[
 Z_y\longrightarrow Z_y^{Z_y}
\]
is a weak equivalence.  Its surjectivity on components makes $1_{Z_y}$
homotopic to a constant map.  Thus $Z_y$ is weakly contractible.  All fibres
of the Kan fibration $p$ are now nonempty and weakly contractible, so
$p\in TF$.  B.2 proves $TC_0=C\cap V$.

Every allowed aspherical interval generator lies in $C\cap V$, while
$\mathcal I_0$ is among them.  Thus
$TC_0\subseteq TC_{\mathrm{all}}\subseteq C\cap V$ proves the first
equality, and right orthogonals give the equality for $F$.  Applying this
to the aspherical core proves the statement for arbitrary non-total
accessible $\mathcal W$.  Finally if $C\cap V=C\cap W$, factor any
$W$-map as a monomorphism followed by an ordinary trivial Kan fibration.
Two out of three makes the monomorphism a $W$-map, hence a $V$-map; the
trivial fibration already lies in $V$.  Therefore $W=V$, and the
test-category correspondence and Theorem 6.1.11 give properness of
$\mathcal W$.  The converse has been proved.

For the total localizer B.3 gives the inclusion in the displayed
emptiness-preserving class.  Use just $(S^0,e)$ for the other inclusion.
Its right lifting class consists of maps with empty source and of maps
with nonempty source that are trivial Kan fibrations, by the disconnected
case above.  Every emptiness-preserving monomorphism lifts against this
class.  The set-generated weak factorization system of B.1 therefore
gives the reverse inclusion.
\end{proof}

\begin{claim}{B.11}{Proved}{A set of intervals can be encoded by one interval}\label{B11}
Under the hypotheses of B.1, let $\{(I_\lambda,e_\lambda)\}_{\lambda\in S}$
be a set of allowed sectioned intervals.  Form their wedge over the
chosen sections,
\[
 I=\bigvee_{\lambda\in S}(I_\lambda,e_\lambda),
 \qquad e:1\longrightarrow I.
\]
Then $I\to1$ is universally $W$-aspherical and this single sectioned
interval generates exactly the same $TC$ and $F$ as the whole set.
If the bunch contains a separating interval, the resulting interval is
separating as well.  In particular the set supplied by B.10 can be encoded
by one $W$-dependent separating interval.
\end{claim}
\begin{proof}
Each chosen section is a monomorphism.  For every $X$, the map
$X\to X\times I$ is built from the $W$-monomorphisms
$X\to X\times I_\lambda$ by pushouts and a transfinite composite.
It belongs to $W$, and its retraction $X\times I\to X$ does too.
Every pointed interval $I_\lambda$ is a retract of $(I,e)$: collapse all
other summands through their maps to $1$.  Consequently every generator
$h(i,I_\lambda,e_\lambda)$ is a retract of $h(i,I,e)$.
Conversely, for a fixed monomorphism $i:K\hookrightarrow L$, attach the
pieces $L\times I_\lambda$ one by one to
$L\cup(K\times I)$.  Each attachment is a pushout of
$h(i,I_\lambda,e_\lambda)$, and their transfinite composite is $h(i,I,e)$.
The two weak saturations, and hence their right orthogonals, coincide.
The inclusions of summands into the wedge are monomorphisms, so two
disjoint endpoint sections in a summand remain disjoint in $I$.
For B.10 the two pointed copies of $\Delta^1$ occur as separate summands;
their chosen vertices are identified, and each pointed copy is still a
retract.  Thus both endpoint lifting directions are retained even though
the packed interval has only one chosen section.
\end{proof}

The conclusion of B.11 does not make the fixed Lawvere interval encode an
arbitrary prescribed localizer.  The packed interval depends on the full
generating bunch and therefore on $W$.  In the accessible simplicial case,
B.10 identifies exactly the part of $W$ that such intervals can encode.

\begin{claim}{B.12}{Refuted}{Properness alone does not suffice on a general strict test category}\label{B12}
For every presheaf category $\widehat A$, every weakly saturated class
generated by sectioned intervals preserves the support
\[
 \operatorname{supp}(X)=\{a\in\operatorname{Ob}A:X(a)\ne\varnothing\}.
\]
Consequently the proposed equality can fail even when the fundamental
localizer is the proper, accessible ordinary localizer $\mathcal W_\infty$
and $A$ is a strict test category.

Explicitly, take $A=\Delta\times[1]$.  A presheaf is a map of simplicial
sets $X_1\to X_0$.  The inclusion
\[
 j:(\varnothing\longrightarrow\Delta^0)
       \longrightarrow(\Delta^0\xrightarrow{1}\Delta^0)
\]
is a monomorphism in $(\mathcal W_\infty)_A$ but is not generated by
\emph{any} bunch of sectioned intervals.  It is the Yoneda map
$y(([0],0))\to y(([0],1))=1$.  Thus B.10 is a theorem about simplicial
sets, not a properness criterion for every presheaf presentation.
\end{claim}
\begin{proof}
Fix $a\in A$.  Evaluation at $a$ preserves products, pushouts,
transfinite colimits and retracts.  The source of
$h(i,I,e)(a)$ contains the copy $L(a)\times\{e(a)\}$, and its target is
$L(a)\times I(a)$, where $I(a)$ is nonempty.  Both are therefore empty
exactly when $L(a)$ is empty.  The set maps preserving emptiness form a
class closed under pushouts, transfinite composites and retracts, as in
B.3.  Thus every interval-generated map has identical source and target
support.

We verify the precise weak equivalences in the example.  For
$X=(X_1\to X_0)$, the inclusion
\[
 J_X:\Delta/X_0\longrightarrow (\Delta\times[1])/X
\]
has a right adjoint $R_X$.  It sends a simplex of $X_i$ to its restriction
in $X_0$.  On objects $R_XJ_X=1$, and the arrow $0\to i$ gives the
counit $J_XR_X\to1$.  The adjunction identifies the hom-sets directly:
a morphism from a simplex in the $0$-layer to one in the $i$-layer is
exactly a simplex morphism after restriction to $X_0$.
These adjunctions are natural in $X$.  Hence $J_X$ is an ordinary weak
equivalence and
\[
 B((\Delta\times[1])/X)\simeq B(\Delta/X_0)\simeq |X_0|.
\]
It follows that the induced localizer on this presheaf category consists
exactly of maps that are ordinary weak equivalences at $0$.

The category $\Delta\times[1]$ is a test category by
\citet[Corollary 4.1.22]{Cisinski2006B}, since $\Delta$ is a test category
and $[1]$ is aspherical.  The just-computed weak equivalences are preserved
by finite products, because evaluation at $0$ is product preserving.
Proposition 4.3.2(g) and Definition 4.3.3 of the same source therefore
make this a \emph{strict} test category.  This argument does not assert
that $\Delta$ has categorical finite products.

The displayed $j$ is the identity at $0$, so it is a weak equivalence.
At $([0],1)$ it is $\varnothing\to\{*\}$; its supports differ, excluding
it from $TC$.  There is also a direct right-lifting witness, independent
of any set replacement.  Write $E=(\varnothing\to\Delta^0)$.
The map $j:E\to1$ lifts against every sectioned-interval generator
$h(i,I,e)$.  Indeed a map from the source of that generator into $E$
forces $L_1=\varnothing$, because its source contains $L_1$ in layer
$1$.  Then $(L\times I)_1=\varnothing$ as well.  At layer $0$ there is
a unique map to $E_0=\Delta^0$, and these maps define the lift.
Thus $j\in F\cap W$ for every possible bunch, whereas $j\notin TF$:
the square of $j$ against itself with identity horizontal maps has no
lift, since $E$ has no global section.  In particular $j$ also lies
outside ${}^{\perp}F$, even if that notation is used for a proper-class
bunch instead of the explicit weak saturation.

The data are finitely presented by the two representables
and their Yoneda map.  The program
\texttt{code/B\_support\_counterexample.py} constructs the finite
$[1]$-diagrams, checks their restrictions, the monomorphism and the two
supports, and checks the terminal objects of the finite categories of
elements.  The $\Delta$ factor is constant in these two diagrams and the
adjunction above verifies its role in every degree.

More generally, under B.1, if $TC=C\cap W$, then \emph{every} $W$-map
preserves support.  Factor it as a monomorphism followed by a member of
$TF$.  The first map is in $TC$ by two out of three; the second has a
section and also preserves support.  In particular, on a test category
with its induced fundamental localizer, every representable collapse
$y(a)\to1$ is a weak equivalence.  Successful interval recognition would
therefore force $\operatorname{Hom}_A(b,a)\ne\varnothing$ for every
pair of objects $a,b$.  This necessary condition fails in the example.
\end{proof}

\begin{claim}{B.13}{Refuted}{The componentwise product case fails even for ordinary weak equivalences}\label{B13}
In case \textup{(b)} of cote 105/020, the manuscript explicitly defines
$W=\underline W_{A/A_0}$ \emph{componentwise} through
$\underline W_T$, under
$\widehat A\simeq\operatorname{Fun}(A_0^{\mathrm{op}},\widehat T)$.
Even for the ordinary proper accessible fundamental localizer
$\underline W=\mathcal W_\infty$, $T=\Delta$, and a finite base $A_0$,
the class generated by \emph{all} permitted sectioned intervals can
be strictly smaller than $C\cap W$.

Take $A_0=BC_2$.  Identify its simplicial presheaves with simplicial
sets carrying a $C_2$-action.  Let $\mathcal E$ be the chaotic category
on the two-element set $C_2$, with its regular action on objects and
the induced action on its four arrows.  Let $\mathcal D=\mathcal
E^{\triangleright}$ be its terminal cone, with fixed cone vertex.
Then
\[
 j:N\mathcal E\hookrightarrow N\mathcal D
 \quad\text{lies in }C\cap W\text{ but not in }TC.
 \tag{B.13.1}
\]
In fact $p:N\mathcal E\to1$ is in $F\cap W$ for \emph{every}
bunch of sectioned intervals and is not in $TF$.  Thus the example
also excludes recognition when $TC$ is interpreted as
${}^{\perp}F$ for a proper-class bunch.
\end{claim}
\begin{proof}
Put $G=C_2$ and $E=N\mathcal E$.  In degree $n$,
$E_n=G^{n+1}$ with the free diagonal action.  Since $\mathcal E$ is
a contractible groupoid, $E$ is Kan and weakly contractible.  Therefore
the underlying map $E\to\Delta^0$ is a trivial Kan fibration.
The cone category has three objects and seven arrows, including its
fixed terminal object.  Both nerves are contractible, so $j$ is an
objectwise ordinary weak equivalence; it is a monomorphism because
the functor is the inclusion of a full subcategory.

We use the elementary equivariant lifting fact that $E\to1$ lifts
against every monomorphism $U\hookrightarrow V$ of degreewise free
$G$-simplicial sets.  Filter $V$ relative to $U$ by its nondegenerate
simplices.  Each new orbit has trivial stabilizer, and hence is
attached by a copy of
\[
 G\times\partial\Delta^n\hookrightarrow G\times\Delta^n.
\]
Choose one orbit representative, lift its boundary by the underlying
trivial Kan fibration, and extend the lift equivariantly.  Freeness
makes the extension well-defined.  Dimension induction and the
colimit give the asserted lift, including dimension zero.

Now consider an arbitrary lifting square for
$h(i,I,e):L\cup(K\times I)\hookrightarrow L\times I$
against $p$.  Here $e:1\to I$ is a global equivariant section.
The top map restricts to an equivariant map $L\times\{e\}\to E$.
It follows that $L$ is degreewise free: every element's stabilizer
must fix its image in the free $G$-set $E_n$.  Consequently
$L\times I$ and its indicated subobject are degreewise free.
The preceding lifting fact supplies the lift.  This argument does
not even need the asphericity of $I$; it applies to every sectioned
interval and therefore to every permitted bunch.

There is no equivariant map $N\mathcal D\to E$, because its fixed
cone vertex would have to map to a fixed vertex of $E$ and no such
vertex exists.  Thus the square with left edge $j$, right edge $p$,
top edge $\operatorname{id}_E$ and bottom edge the collapse has no
lift.  Since $p\in F$, this excludes $j$ from both the weak saturation
of the interval generators and ${}^{\perp}F$.  Taking the left edge
$\varnothing\to1$ also shows $p\notin TF$.

All five preliminary hypotheses in cote 105/016--018 hold for this
$W$.  The constant interval $\Delta^1$, with trivial $G$-action,
is separating and its product projection is an objectwise weak
equivalence.  The same is true for the projection of every permitted
objectwise aspherical interval.  Pushouts of weak-equivalence
monomorphisms, filtered colimits and retracts are detected on the
underlying simplicial sets and preserve weak equivalences.
The class is accessible as well.  For clarity, this does not rely on
the failed interval recognition: a countable submap of an underlying
weak equivalence can be enlarged to a countable weak-equivalence
submap by adjoining representatives and nullhomotopies for its
homotopy-group tests after $\operatorname{Ex}^{\infty}$.
Close all added simplices under $G$ and simplicial operators; this
remains countable since $G$ is finite.  Iterate the countably many
tests, including components and all basepoints.  Such countable
equivariant weak-equivalence submaps form an $\aleph_1$-filtered cover
of each map in $W$.  Countable simplicial $G$-sets are
$\aleph_1$-presentable, and weak equivalences are preserved by
filtered colimits.  This proves accessibility directly.

The example can also be placed over a finite \emph{aspherical} base.
Adjoin an initial object $a$ to $BG$.  At the old group object use
$E\hookrightarrow N\mathcal D$; at $a$ use
$\Delta^0\xrightarrow{1}\Delta^0$, with the unique restriction maps
to that value.  This is a componentwise weak-equivalence
monomorphism.  The same lifting proof works at the group object,
and the point component has the unique compatible lift.  A global
section would still require a fixed vertex at the group object.
Thus the same failure occurs for this aspherical two-object,
four-arrow base.
\end{proof}

The finite certificate \texttt{code/B\_equivariant\_counterexample.py}
constructs these categories and actions, verifies their compositions,
terminal objects and fixed vertices, and checks the nerve actions and
simplicial operators in bounded degrees.  The proof supplies the
unbounded Kan and lifting statements.  Representable horn inclusions
do not contradict B.13: on $BG$ they attach free $G$-cells and generate
the projective cofibrations.  They do not generate all objectwise
weak-equivalence monomorphisms when cofibrations are prescribed to be
\emph{all} monomorphisms.  The proposed positive inference from
objectwise horns to the manuscript's $C$ was considered and rejected
by this explicit example.

In particular B.12 concerns the absolute localizer in case (a); its
evaluation-at-$0$ weak equivalences must not be substituted for the
componentwise class of case (b).  B.13 answers the latter case with its
actual definition, even for $\mathcal W_\infty$.  The special case
$A_0=e$, $T=\Delta$ retains the exact positive theorem B.10.
